\documentclass[11pt]{article}
\usepackage{color}
%\input{rgb}

%----------Packages----------
\usepackage{amsmath}
\usepackage{amssymb}
\usepackage{amsthm}
%\usepackage{amsrefs}
\usepackage{dsfont}
\usepackage{mathrsfs}
\usepackage{stmaryrd}
\usepackage[all]{xy}
\usepackage[mathcal]{eucal}
\usepackage{verbatim}  %%includes comment environment
\usepackage{fullpage}  %%smaller margins
%----------Commands----------

%%penalizes orphans
\clubpenalty=9999
\widowpenalty=9999





%% bold math capitals
\newcommand{\bA}{\mathbf{A}}
\newcommand{\bB}{\mathbf{B}}
\newcommand{\bC}{\mathbf{C}}
\newcommand{\bD}{\mathbf{D}}
\newcommand{\bE}{\mathbf{E}}
\newcommand{\bF}{\mathbf{F}}
\newcommand{\bG}{\mathbf{G}}
\newcommand{\bH}{\mathbf{H}}
\newcommand{\bI}{\mathbf{I}}
\newcommand{\bJ}{\mathbf{J}}
\newcommand{\bK}{\mathbf{K}}
\newcommand{\bL}{\mathbf{L}}
\newcommand{\bM}{\mathbf{M}}
\newcommand{\bN}{\mathbf{N}}
\newcommand{\bO}{\mathbf{O}}
\newcommand{\bP}{\mathbf{P}}
\newcommand{\bQ}{\mathbf{Q}}
\newcommand{\bR}{\mathbf{R}}
\newcommand{\bS}{\mathbf{S}}
\newcommand{\bT}{\mathbf{T}}
\newcommand{\bU}{\mathbf{U}}
\newcommand{\bV}{\mathbf{V}}
\newcommand{\bW}{\mathbf{W}}
\newcommand{\bX}{\mathbf{X}}
\newcommand{\bY}{\mathbf{Y}}
\newcommand{\bZ}{\mathbf{Z}}

%% blackboard bold math capitals
\newcommand{\bbA}{\mathbb{A}}
\newcommand{\bbB}{\mathbb{B}}
\newcommand{\bbC}{\mathbb{C}}
\newcommand{\bbD}{\mathbb{D}}
\newcommand{\bbE}{\mathbb{E}}
\newcommand{\bbF}{\mathbb{F}}
\newcommand{\bbG}{\mathbb{G}}
\newcommand{\bbH}{\mathbb{H}}
\newcommand{\bbI}{\mathbb{I}}
\newcommand{\bbJ}{\mathbb{J}}
\newcommand{\bbK}{\mathbb{K}}
\newcommand{\bbL}{\mathbb{L}}
\newcommand{\bbM}{\mathbb{M}}
\newcommand{\bbN}{\mathbb{N}}
\newcommand{\bbO}{\mathbb{O}}
\newcommand{\bbP}{\mathbb{P}}
\newcommand{\bbQ}{\mathbb{Q}}
\newcommand{\bbR}{\mathbb{R}}
\newcommand{\bbS}{\mathbb{S}}
\newcommand{\bbT}{\mathbb{T}}
\newcommand{\bbU}{\mathbb{U}}
\newcommand{\bbV}{\mathbb{V}}
\newcommand{\bbW}{\mathbb{W}}
\newcommand{\bbX}{\mathbb{X}}
\newcommand{\bbY}{\mathbb{Y}}
\newcommand{\bbZ}{\mathbb{Z}}

%% script math capitals
\newcommand{\sA}{\mathscr{A}}
\newcommand{\sB}{\mathscr{B}}
\newcommand{\sC}{\mathscr{C}}
\newcommand{\sD}{\mathscr{D}}
\newcommand{\sE}{\mathscr{E}}
\newcommand{\sF}{\mathscr{F}}
\newcommand{\sG}{\mathscr{G}}
\newcommand{\sH}{\mathscr{H}}
\newcommand{\sI}{\mathscr{I}}
\newcommand{\sJ}{\mathscr{J}}
\newcommand{\sK}{\mathscr{K}}
\newcommand{\sL}{\mathscr{L}}
\newcommand{\sM}{\mathscr{M}}
\newcommand{\sN}{\mathscr{N}}
\newcommand{\sO}{\mathscr{O}}
\newcommand{\sP}{\mathscr{P}}
\newcommand{\sQ}{\mathscr{Q}}
\newcommand{\sR}{\mathscr{R}}
\newcommand{\sS}{\mathscr{S}}
\newcommand{\sT}{\mathscr{T}}
\newcommand{\sU}{\mathscr{U}}
\newcommand{\sV}{\mathscr{V}}
\newcommand{\sW}{\mathscr{W}}
\newcommand{\sX}{\mathscr{X}}
\newcommand{\sY}{\mathscr{Y}}
\newcommand{\sZ}{\mathscr{Z}}

\newcommand{\id}{\mathsf{id}}
\newcommand{\Char}{\mathrm{char}}

\renewcommand{\phi}{\varphi}

%\renewcommand{\emptyset}{\O}

\providecommand{\abs}[1]{\lvert #1 \rvert}
\providecommand{\norm}[1]{\lVert #1 \rVert}


\providecommand{\x}{\times}




\providecommand{\ar}{\rightarrow}
\providecommand{\arr}{\longrightarrow}



\newcommand{\head}[1]{
	\begin{center}
		{\large #1}
		\vspace{.2 in}
	\end{center}
	
	\bigskip 
}



%----------Theorems----------

\newtheorem{theorem}{Theorem}[section]
\newtheorem{proposition}[theorem]{Proposition}
\newtheorem{lemma}[theorem]{Lemma}
\newtheorem{corollary}[theorem]{Corollary}

\theoremstyle{definition}
\newtheorem{definition}[theorem]{Definition}
\newtheorem*{definition*}{Definition}
\newtheorem{nondefinition}[theorem]{Non-Definition}
\newtheorem{exercise}[theorem]{Exercise}



\numberwithin{equation}{subsection}


%----------Title-------------
\newcommand{\hide}[1]{{\color{red} #1}} 
\newcommand{\com}[1]{{\color{blue} #1}} 
\newcommand{\meta}[1]{{\color{green} #1}} 

\begin{document}

\pagestyle{plain}


%%---  sheet number for theorem counter
\setcounter{section}{1}   

\head{Abstract Linear Algebra, Sections 41\&51, Autumn 2025\\ HOMEWORK 5} 

This homework is due to {\bf November 10}, midnight. The total grade is 30 points. Please upload your solution on Canvas/Gradescope.

\begin{enumerate}

\item (\emph {3 points}) Let $\bbF$ be a field. Recall from Homework 2 that the \emph{characteristic} of $\bbF$, often denoted $\Char(\bbF)$, is the smallest positive (natural) number $n$ such that
$$n\cdot 1 = \underbrace{1+\ldots +1}_{\text{$n$ times}} =0 $$
if such number $n$ exists, and $0$ otherwise. Suppose that the field $\bbF$ is finite. Show that $|\bbF|=p^n$, where $p=\Char(\bbF)$ and $n$ is a natural number.
\item[] Shown in Homework 2: By definition of the characteristic, the resulting subring generated by $1\in \bbF$ is isomorphic to $\bbZ/p\bbZ$ where $p$ is a prime number. Similarly, the prime field within $\bbF$ is isomorphic to $\bbF_p$, the finite field with $p$ elements. Consider $\bbF$ a vector space over $\bbF_p$. Since $\bbF$ is finite, the vector space has a finite dimension $n$ (a natural number). Then, a finite dimensional vector space of $n$ over $\bbF_p$ has exactly $p^n$ elements. Therefore $|\bbF|=p^n$.

\item Let $A \in M_{m\times n}(\bbF)$ be a matrix. We say that a matrix $B\in M_{n\times m}(\bbF)$ is a \emph{left inverse} if $B\cdot A= I_{n}$ and a matrix $C \in M_{n\times m}(\bbF)$ is a \emph{right inverse} if $A\cdot C = I_m$.

\begin{enumerate}
\item (\emph {1 point}) Suppose that a square matrix $A\in M_{n\times n}(\bbF)$ has both a left inverse $B \in M_{n\times n}(\bbF)$ and a right inverse $C \in M_{n\times n}(\bbF)$. Show that $B=C$, and so, $A$ is invertible.
\item[] By computation: $B=B\cdot I_n=B\cdot (A\cdot C)=(B\cdot A)\cdot C=I_n\cdot C=C$. Thus $B=C$. Since there exists a matrix $B$ such that A has a two-sided inverse, A is invertible.
\item (\emph {2 points}) Let $A=\begin{bmatrix} 1 & -1\\ 0 & 1\\ 1 & 1\end{bmatrix}$. Show that $A$ is left invertible (i.e. $A$ has a left inverse).
\item[] Deconstruct $A$ into column vectors in $\bbF^3$: $c_1=\begin{bmatrix}1\\0\\1\end{bmatrix}$, $c_2=\begin{bmatrix}-1\\1\\1\end{bmatrix}$. A left inverse of $A$, $B$, must be a $2\times 3$ matrix such that $AB=I_2$. Construct $B$ as $\begin{bmatrix}a&b&c\\d&e&f\end{bmatrix}$. Further, $BA$ must have columns such that $Bc_1=e_1=\begin{bmatrix}1\\0\end{bmatrix}$, and $Bc_2=e_2=\begin{bmatrix}0\\1\end{bmatrix}$. Expanding below:
\item[] For $Bc_1$: $\begin{cases}
    a\cdot 1+b\cdot 0+c\cdot 1=1 \\ d\cdot 1+e\cdot 0+f\cdot 1=0
\end{cases}$. For $Bc_2$: $\begin{cases}
    a\cdot-1+b\cdot 1+c\cdot 1=0 \\ d\cdot -1+e\cdot1+f\cdot 1=1
\end{cases}$. 
\item[] Rows (1) and (3): $\begin{cases}a+c=1\\-a+b+c=0\end{cases}$. Rows (2) and (4): $\begin{cases}d+f=0\\-d+e+f=1\end{cases}$.
\item[] Solving the system for constants: $c=1-a \Rightarrow b=2a-1$, so top row $= (a, 2a-1, a-1)$. Similarly, $f=-d \Rightarrow e=1+2d$, so bottom row $= (d, 1+2d, -d)$.
\item[] Then, the general left inverse is given by $B(a,d)=\begin{bmatrix}a&2a-1&1-a\\d&1+2d&-d\end{bmatrix}$, for $a,d\in \bbF$. Because there exists explicit $a,d$ such that $BA=I_2$, $A$ is left-invertible.
\item (\emph {2 points}) Let $A=\begin{bmatrix} 1 & -1\\ 0 & 1\\ 1 & 1\end{bmatrix}$. Show that $A$ has at least two left inverses.
\item[] From the general form given above, we can assign two different pairs $(a,d)$:
\item[] $B_1=B(0,0)=\begin{bmatrix}0&-1&-1\\0&1&0\end{bmatrix}$, where $B_1c_1=e_1$ and $B_1c_2=e_2$ holds.
\item[] $B_2=B(1,0)=\begin{bmatrix}1&1&0\\0&1&0\end{bmatrix}$, where $B_2c_1=e_1$ and $B_2c_2=e_2$ holds.
\item[] $B_1$ and $B_2$ are clearly different matrices that hold $BA=I_2$, thus $A$ has at least two distinct left inverses. Further, since $a,d$ are arbitrary to the general form, there are infinitely many left inverses available over infinite fields.
\end{enumerate}

\item A square matrix $A \in M_{n\times n}(\bbF)$ is \textit{nilpotent} if there exists an integer $n \geq 0$ such that $A^n = 0$. 

\begin{enumerate}
\item (\emph{4 points}) Show that a matrix $A$ of the form
$$A = \begin{bmatrix} 
0&a_{12} & a_{13} & \ldots & a_{1n} \\
0& 0& a_{23}& \ldots & a_{2n} \\
\vdots& & \ddots & &  \vdots \\
0& \ldots &\ldots & 0 & a_{(n-1)n} \\
0& \ldots & \ldots & 0 & 0
\end{bmatrix}  \in M_{n\times n}(\bbF)$$
is nilpotent. (The matrix of this form is called \emph{strictly upper triangular}.)
\item[] Here, $A=(a_{ij})=0$ satisfies for all $i\geq j$, defining an upper triangular matrix. Proof: For matrix powers, the $(i,j)$-entry of $A^k$ is the sum of products in the form $a_{i,i_1},a_{i_1,i_2},\ldots,a_{i_{k-1},j}$, where indices create a chain $i<i_1<i_2<\ldots<j$ whenever a factor is 0 (because $a_{pq}\neq 0$ only when $p<q$ in an upper triangular matrix). Thus, for non-zero products contributing to $(A^k)_{ij}$, we must have $i<j$ and strictly increasing intermediate indices of length $\geq k$. However, for an $n\times n$ matrix, the maximum increase in indices from $i$ to $j$ is at most $n-1$. Any chain of $k$ multiplications requires $i-j\geq k$. But if $k\geq n$, there is no possible chain of length $k$ (because there are only $n$ indices), so all entries of $A^n$ are zero. Thus, $A^n=0$ and A in the given form is nilpotent.

\item (\emph{3 points}) Let $A\in M_{n\times n}(\bbF)$ be a nilpotent matrix. Show that the matrix $I-A$ is invertible. Here $I$ is the diagonal matrix with $1$'s on the diagonal. {\it Hint: Recall that $1+x+x^2+\ldots + x^m=\frac{1-x^{m+1}}{1-x}$}.
\item[] Following the finite geometric sum, define $S=I+A+A^2+\ldots +A^{m-1}$. Using distributivity and the associativity of matrix multiplication, compute the product $(I-A)S$:
\item[] $(I-A)S=S-AS=(I+A+A^2+\ldots+A^{m-1})-(A+A^2+\ldots +A^m)=I-A^m$.
\item[] Because A is nilpotent, $A^m=0$, thus $(I-A)S=I-0=I$. Now compute $S(I-A)$: 
\item[] $S(I-A)=S-SA=(I+A+A^2+\ldots+A^{m-1})-(A+A^2+\ldots+A^m)=I-A^m$.
\item[] This computation also yields $I$, since commutativity holds under powers of equal base. Therefore, $S$ is both a left and right inverse for $I-A$, proving $I-A$ is invertible.
\end{enumerate}


\item Let $k\in\mathbb{R}$ be a constant and consider the system of equations
$$\begin{cases} x+ky-2z=0\\ kx+y+z=0\\ 4x+ky-8z=0\end{cases}.$$
Use row reduction to determine the answers to the following:

\begin{enumerate}

\item (\emph{2 points}) For which values of $k$ does this system have \emph{only} the solution $x=0$, $y=0$, $z=0$?
\item[] The coefficient matrix is given by $M=\begin{bmatrix}1&k&-2\\k&1&1\\4&k&-8\end{bmatrix}$. The homogeneous system only contains the trivial solution when det$(M)\neq 0$. Thus, computing the determinant
\item[] det$(M)=1\cdot \det\begin{bmatrix}1&1\\k&-8\end{bmatrix}-k\cdot\det\begin{bmatrix}k&-2\\k&-8\end{bmatrix}+4\cdot\det\begin{bmatrix}k&-2\\1&1\end{bmatrix}$
\item[] $\det(M)=1(1\cdot-8-k\cdot1)-k(k\cdot-8-k\cdot-2)+4(k\cdot1-1\cdot-2)=(-8-k)-k(-8k+2k)+4(k+2)$
\item[] $\Rightarrow (-8-k)+(8k^2-2k^2)+(4k+8)=6k^2+3k=3k(2k+1)=\det(M)$.
\item[] Solving for $k$ such that $\det(M)=3k(2k+1)\neq0$ yields $k=-\frac{1}{2},0$.

\item (\emph{3 points}) In the case(s) where there is more than just the solution $x=0$, $y=0$, $z=0$, describe all possible solutions.
\item[] Case when $k=0$: $\begin{bmatrix}1&0&-2\\0&1&1\\4&0&-8\end{bmatrix}\Rightarrow\begin{bmatrix}1&0&-2\\0&1&1\\0&0&0\end{bmatrix}\Rightarrow x_2=-x_3,\; x_1=2x_3\Rightarrow x_3\begin{bmatrix}2\\-1\\1\end{bmatrix}$.
\item[] Thus, all possible non-trivial solutions when $k=0$ exist in the form $t(2,-1,1)$ for $t\in\bbR$.
\item[] Case when $k=-\frac{1}{2}$: $\begin{bmatrix}1&-\frac{1}{2}&-2\\-\frac{1}{2}&1&1\\4&-\frac{1}{2}&-8\end{bmatrix}\Rightarrow\begin{bmatrix}1&-\frac{1}{2}&-2\\0&\frac{3}{4}&0\\0&0&0\end{bmatrix}\Rightarrow x_2=0, \; x_1=2x_3\Rightarrow x_3\begin{bmatrix}2\\0\\1\end{bmatrix}$.
\item[] Thus, all possible non-trivial solutions when $k=-\frac{1}{2}$ exist in the form $t(2,0,1)$ for $t\in\bbR$.

\item (\emph{2 points}) Use your answers to the above to determine the values of $k\in\mathbb{R}$ such that the vectors
$$\begin{bmatrix} 1\\ k\\ 4\end{bmatrix}, \quad \begin{bmatrix} k\\ 1\\ k\end{bmatrix}, \quad \begin{bmatrix} -2\\ 1\\ -8\end{bmatrix}$$
are linearly independent vectors in $\mathbb{R}^3$.
\item[] Vectors are linearly independent if the matrix formed has a non-zero determinant. Therefore, as shown above, the three given vectors are independent when $k\neq0$ and $k\neq -\frac{1}{2}$.

\end{enumerate}

\item Use row reduction to solve the following systems of linear equations over the real numbers $\bbR$. If a system has infinitely many solutions, then describe all possible solutions in terms of free variables. If a system has no solution, then show this fact.

\begin{enumerate}

\item (\emph{4 points}) $\begin{cases} x+y-2z=-3\\ 2x-y+3z=7\\ x-2y+5z=1\end{cases}$
\item[] $M=\left[\begin{array}{ccc|c}1&1&-2&-3\\2&-1&3&7\\1&-2&5&1\end{array}\right]\Rightarrow\left[\begin{array}{ccc|c}1&1&-2&-3\\0&-3&7&13\\0&-3&7&4\end{array}\right]\Rightarrow\left[\begin{array}{ccc|c}1&1&-2&-3\\0&-3&7&13\\0&0&0&-9\end{array}\right]$.
\item[] The system has no solutions due to contradiction. $0\neq-9$ is inconsistent.

\item (\emph{4 points}) $\begin{cases} x+2y=1\\ 2x+5y=2\\ 3x+6y=3\end{cases}$
\item[] $M=\left[\begin{array}{cc|c}1&2&1\\2&5&2\\3&6&3\end{array}\right]\Rightarrow\left[\begin{array}{cc|c}1&2&1\\0&1&0\\0&0&0\end{array}\right]$.
\item[] Here, the third row provides no constraint ($n$ exceeds $m$ with no contradiction), thus the system is consistent and has a unique solution where $x_1=1$ and $x_2=0$.

\end{enumerate}

\end{enumerate}

\end{document}
